\begin{abstract}
A compreensão de orientações técnicas de saúde pode ser dificultada pelo uso de vocabulário especializado, frases longas e baixa adequação ao perfil comunicacional de usuários finais em aplicações mHealth. Este trabalho tem como objetivo investigar, implementar e avaliar uma solução baseada em Modelos de Linguagem de Grande Porte para simplificação e adaptação comunicacional de textos técnicos de saúde, preservando o conteúdo essencial do material original. Para isso, foi desenvolvida a SimplyHealth, contextualizada na Takere, uma plataforma no-code para aplicações mHealth baseadas em planos de cuidado, tratada como ambiente de referência sem integração plena. A solução recebe textos livres ou arquivos PDF com orientações de saúde, considera perfis simulados de pacientes e gera versões simplificadas, glossário, perguntas frequentes e chat contextual. A adaptação é estritamente comunicacional, não clínica, utilizando aspectos como idade, escolaridade e área de formação apenas para ajustar vocabulário, tom, nível de explicação e forma de apresentação. A implementação inclui prompts com regras de preservação do conteúdo, guardrails em duas camadas — verificação determinística e avaliação por LLM — para mitigar alterações de sentido, e o Índice de Legibilidade de Flesch adaptado ao português (Flesch-PT) como indicador automático de legibilidade. A avaliação funcional e linguística utilizou dois materiais de saúde, três perfis simulados e 18 execuções principais. Os guardrails atuaram em 8 das 18 execuções, bloqueando simplificações com alterações identificadas como problemáticas. Os resultados de legibilidade mostraram aumento do Flesch-PT em 15 dos 18 casos, com média de 46,8 nos textos originais e 58,2 nas versões simplificadas, resultando em ganho médio de 11,4 pontos. Os achados indicam viabilidade funcional e linguística da abordagem no escopo avaliado, sem representar validação clínica ou comprovação de compreensão por pacientes reais.
\end{abstract}
